% ============================================================================
\chapter{Ступень III. Квантовая механика, электродинамика Максвелла и демоны Лапласа и Максвелла}
\label{ch:course-stage3-quantum}
% ============================================================================

\section{Урок 8. Мёртвая зона затвора \texorpdfstring{$\Delta\varphi_{\mathrm{disc}} = 41$}{dphi = 41} и полное квантование всех наблюдаемых}

Почему атомы не излучают энергию непрерывно и не падают на ядро? Посмотрим на то, как 13-узловой кластер кубооктаэдра ($N_\varphi = 1 + 12 = 13$) делит полное фазовое кольцо ($N_{\mathrm{ring}} = 512$) с учётом 8 октантов ($|{\triangle}| = 8$):
\begin{equation}
  \label{eq:course-heisenberg-41}
  \Delta\varphi_{\mathrm{disc}}(3)
  \;=\; \frac{\phaseRingCount(3) + N_\varphi(3) + |{\triangle}(3)|}{N_\varphi(3)}
  \;=\; \frac{512 + 13 + 8}{13}
  \;=\; \frac{533}{13}
  \;=\; \mathbf{41\text{ тик}}.
\end{equation}
Число $533 = 13 \times 41$ делится на $13$ \textbf{нацело} исключительно при $d = 3$!

Любой фазовый импульс ворот $|\Phi| < 41$ попадает внутрь \textbf{мёртвой зоны затвора}, где целочисленное деление $\hat{n}_E(\Phi) = \lfloor|\Phi|/41\rfloor = \mathbf{0}$ даёт тождественный ноль излучённых квантов энергии. А поскольку сверху фазовый импульс ограничен полупериодом противофазы $|\Phi|\le\Phi_{\mathrm{Pauli}}(3) = 256$, максимальное число квантов энергии, которое клетка может испустить или поглотить за один такт, равно числу квадратных граней кубооктаэдра:
\[
  \hat{n}_E(\Phi) \;=\; \bigl\lfloor |\Phi| / 41 \bigr\rfloor \;\in\; \{0, 1, 2, 3, 4, 5, 6 = |\square(3)|\},
  \qquad
  256 \;=\; 6 \times 41 + 10 \quad \bigl(10 = Z(\mathrm{Ne})\bigr).
\]
В результате в автомате \textbf{не существует ни одной непрерывной наблюдаемой}: однотактовый спинорный регистр клетки весит ровно $B_1 = 4\times 9 = 36 = 6^2$ бит (по числу рёбер $\mathcal{C}_{13}$), а полное двухтактовое динамическое состояние клетки $(P,Q)$ весит ровно:
\[
  \mathbf{B_2(3) \;=\; 2\,B_1(3) \;=\; 8 \times 9 \;=\; 6 \times 12 \;=\; 72\text{ бита}}.
\]

\section{Урок 9. Корпускулярно-волновой дуализм и вывод правила Борна \texorpdfstring{$P(\mu) = \mu^2 / 512$}{P(mu) = mu\^2 / 512} из пина переноса 9-битного сумматора}

В учебниках XX~века корпускулярно-волновой дуализм подавали как мистику: «электрон --- это то волна, то частица, а квадрат модуля амплитуды $|\psi|^2$ задаёт вероятность щелчка детектора по постулату Борна».

А теперь посмотрим, как работает деление разности фаз на $512$ с остатком в каждом такте:
\[
  \Delta\varphi \;=\; 512\,q + r, \qquad q\in\{-1,0,+1\},\quad r\in(-256,256].
\]
\begin{itemize}[noitemsep]
\item \textbf{Волна} --- это 9-битный остаток $r$: пока $|\Delta\varphi|\le 256$, частное $q = 0$, и остаток $r$ свободно бежит и интерферирует по $12$ рёбрам кубооктаэдра.
\item \textbf{Корпускула (частица)} --- это целое частное $q$ (бит переноса \texttt{carry\_out} за границу кольца $512$): целое число $1$ по законам целочисленной арифметики не может разделиться пополам на полупрозрачном зеркале и всегда срабатывает как неделимый щелчок $q = \pm 1$ в одном канале!
\end{itemize}

А откуда берётся закон Борна $P(\mu) = \mu^2 / 512$ без единого вероятностного постулата? Поскольку момент голономии звезды $z(\mathbf{y})z^*(\mathbf{x})$ билинеен по амплитудам ячеек, волна с целочисленной амплитудой $\mu\in[0,16]$ приносит на атом детектора квадратичную фазовую энергию $\Phi_{\mathrm{wave}} = \mu^2\in[0,256]$. Внутренняя фаза самого атома детектора $\varphi_{\mathrm{bg}}\in\{0,1,\dots,511\}$ макронаблюдателю неизвестна и за полный цикл равномерно пробегает все $512$ точек кольца $\mathbb{Z}_{512}$. Подадим $\varphi_{\mathrm{bg}}$ и $\mu^2$ на обычный 9-битный сумматор:
\[
  \texttt{\{carry\_click, sum\_wave[8:0]\}} \;=\; \varphi_{\mathrm{bg}} + \mu^2.
\]
В скольких в точности из $512$ значений фона $\varphi_{\mathrm{bg}}\in\{0,\dots,511\}$ сумма $\varphi_{\mathrm{bg}} + \mu^2 \ge 512$ перешагнёт через верхний край кольца и выдаст на пине переноса корпускулярный щелчок \texttt{carry\_click = 1}? Неравенство $\varphi_{\mathrm{bg}} + \mu^2 \ge 512$ выполняется тогда и только тогда, когда $\varphi_{\mathrm{bg}} \in \{512 - \mu^2,\ \dots,\ 511\}$, а в этом интервале содержится ровно $(511 - (512 - \mu^2) + 1) = \mathbf{\mu^2}$ значений фона из $512$:
\begin{equation}
  \label{eq:course-born}
  N_{\mathrm{clicks}}(\mu) \;\equiv\; \mu^2
  \qquad\Longrightarrow\qquad
  P(\mu) \;=\; \frac{\mu^2}{512} \;=\; \frac{|\psi|^2}{\phaseRingCount(3)}!
\end{equation}
Правило Борна --- это буквально частота срабатывания пина переноса \texttt{carry\_out} 9-битного сумматора!

\begin{quote}
\small\textbf{Историческая справка: Макс Борн (1926) и Вернер Гейзенберг (1927).}
Летом 1926~года Макс Борн в статье о квантовой теории столкновений добавил в сноске знаменитую поправку: вероятность обнаружить частицу пропорциональна не самой волновой амплитуде $\psi$, а её квадрату $|\psi|^2$. В стандартной непрерывной квантовой механике это правило так и осталось независимым постулатом («постулатом измерения»), не выводимым из уравнения Шрёдингера. Годом позже, в 1927~году, Вернер Гейзенберг сформулировал принцип неопределённости $\Delta x\,\Delta p \ge \hbar/2$. На дискретном носителе оба этих принципа перестают быть постулатами: порог дискретности $\Delta\varphi_{\mathrm{disc}} = 41$ ($r_{\mathcal{M}} = 41/512 \leftrightarrow 1/(4\pi)$) --- это минимальный шаг срабатывания делителя кольца на 13-узловом ядре, а квадратичный закон Борна $\mu^2/512$ --- число переполнений 9-битного сумматора за полный период фоновой фазы.
\end{quote}

\section{Урок 10. Теорема Стокса--Гаусса--Черна (\texorpdfstring{$d^2 \equiv 0$}{d\^2 = 0}) и дискретные уравнения Максвелла}

Обойдём теперь квадратную грань кубооктаэдра $f = (x_0\to x_1\to x_2\to x_3\to x_0)$ по четырём ориентированным рёбрам. На каждом ребре разность узловых фаз раскладывается на целое число полных оборотов и остаток: $(d_0\varphi)_e = 512\,k_e + A_e$.

Сложим четыре разности вокруг замкнутого квадрата:
\[
  (d_1 d_0\varphi)_f \;=\; (\varphi_1 - \varphi_0) + (\varphi_2 - \varphi_1) + (\varphi_3 - \varphi_2) + (\varphi_0 - \varphi_3) \;\equiv\; \mathbf{0}!
\]
Но поскольку $(d_0\varphi)_e = 512\,k_e + A_e$, из равенства нулю полной суммы мгновенно следует \textbf{дискретная теорема Стокса--Черна}:
\begin{equation}
  \label{eq:course-stokes-chern}
  B_f \;:=\; (d_1 A)_f \;=\; A_{01} + A_{12} + A_{23} + A_{30}
  \;=\; -512\,(k_{01} + k_{12} + k_{23} + k_{30})
  \;=\; \mathbf{512\,n_f \;\in\; 512\,\mathbb{Z}}!
\end{equation}
Магнитный поток $B_f$ через любую плакету \textbf{строго квантован} кратными $512$, где $n_f = -(d_1 k)_f\in\mathbb{Z}$ --- топологическое вихревое число Черна! А если сложить потоки $B_f$ по всем $6$~граням замкнутого куба, каждое из $12$~рёбер куба входит в две соседние грани с противоположными знаками и сокращается:
\[
  (d_2 B)_{\mathrm{cube}} \;=\; (d_2\circ d_1 A)_{\mathrm{cube}} \;\equiv\; \mathbf{0} \qquad (\nabla\cdot\mathbf{B} = 0\text{: магнитных монополей не существует!}),
\]
а разность по времени даёт точный закон электромагнитной индукции Фарадея $\Delta_t B_f + (d_1 E)_f = 512\,\Delta_t n_f$.

\begin{quote}
\small\textbf{Историческая справка: Джеймс Клерк Максвелл (1865), Анри Пуанкаре (1895), Жорж де Рам (1931) и Шиинг-Шен Черн (1946).}
В 1895~году Анри Пуанкаре заметил фундаментальный геометрический факт: \emph{«граница границы любого тела всегда пуста»} ($\partial^2 \equiv 0$, или в двойственных разностях коцепей Жоржа де Рама $d^2 = d_{k+1}\circ d_k \equiv 0$): у отрезка граница состоит из двух концов, у квадрата граница состоит из четырёх отрезков-сторон, чьи концы попарно смыкаются и взаимно уничтожаются, а у куба граница состоит из шести квадратных граней, чьи двенадцать рёбер попарно проходятся в противоположных направлениях и тоже сокращаются в точный ноль. Когда Джеймс Клерк Максвелл в 1865~году записывал уравнения электродинамики ($\nabla\cdot\mathbf{B} = 0$ и $\nabla\times\mathbf{E} + \partial_t\mathbf{B} = 0$), а Шиинг-Шен Черн в 1946~году вводил целочисленные топологические классы расслоений, они описывали на языке дифференциальных форм именно это тождество телескопического сокращения разностей вокруг грани и куба.
\end{quote}

\medskip
\noindent\textbf{Что такое электрический ток на самом деле: почему по проводу никто никуда не бежит, а передаётся эстафета фазы $\mathcal{S}_3(\Delta\varphi)$!}
В школьных учебниках электрический ток рисуют как «направленное движение заряженных шариков-электронов, бегущих по проводу от батарейки к лампочке». Но любой студент, подставивший числа в классическую формулу Друде, обнаруживает абсурд: дрейфовая скорость электронов в медном проводе составляет около $0{,}07\text{ мм/с}$ (чтобы доползти по метровому проводу до люстры, «шарику» нужно четыре часа, а в сети переменного тока $50\text{ Гц}$ он вообще топчется на месте!), тогда как сигнал и энергия приходят к лампочке со скоростью света в линии!

На нашем носителе этот парадокс исчезает мгновенно:
\begin{enumerate}[label=\textbf{\arabic*.},noitemsep]
\item \textbf{Кристаллическая решётка проводника --- это цепочка кубооктаэдров $O_2(3)$:}
  все лучшие металлические проводники --- медь ($\mathrm{Cu},\ Z = 29 = 4\times 7 + 1$), серебро ($\mathrm{Ag},\ Z = 47$), золото ($\mathrm{Au},\ Z = 79$) и алюминий ($\mathrm{Al},\ Z = 13 = N_\varphi(3)$) --- кристаллизуются в гранецентрированную кубическую решётку (ГЦК), у которой первая координационная сфера из $12$ ближайших соседей каждого атома \textbf{в точности совпадает с нашим 12-вершинным кубооктаэдром $O_2(3)$}! Атомы сидят на своих местах и никуда не бегут.
\item \textbf{Напряжение $V = \mathcal{C}_3(\Delta\varphi)$ и ток $I = \mathcal{S}_3(\Delta\varphi)$ как две проекции одной фазы:}
  когда цепь разомкнута, соседние ячейки вдоль провода синфазны ($\Delta\varphi = \varphi_{x+e} - \varphi_x = 0$): чётная проекция ротатора максимальна ($V = \mathcal{C}_3(0) = +128$, напряжение холостого хода), а нечётная проекция равна нулю ($I = \mathcal{S}_3(0) = 0$, тока нет). Замкнуть цепь --- значит создать между соседними ячейками кубооктаэдра сдвиг фазы $\Delta\varphi \in (0, 128]$. Тогда по фундаментальному тождеству грани кубооктаэдра:
  \begin{equation}
    \label{eq:course-ohm-rotator}
    V(\Delta\varphi) + I(\Delta\varphi) \;=\; \mathcal{C}_3(\Delta\varphi) + \mathcal{S}_3(\Delta\varphi) \;=\; (128 - \Delta\varphi) + \Delta\varphi \;\equiv\; \mathbf{128 \;=\; Q_\square(3)}!
  \end{equation}
  \textbf{Электрический ток $I = \mu_x\mu_y\mathcal{S}_3(\Delta\varphi)$ --- это не перенос кусков вещества сквозь металл, а волна (эстафета) фазового кручения $\mathcal{S}_3(\Delta\varphi)$, передающаяся от одного кубооктаэдра к соседнему со скоростью света в линии!}
\item \textbf{Закон Джоуля--Ленца, согласование импедансов ($P_{\max} = 64^2 = 4096$) и сверхпроводимость:}
  передаваемая через звено мощность равна произведению чётного напряжения на нечётный фазовый ток:
  \begin{equation}
    \label{eq:course-power-4096}
    P(\Delta\varphi) \;=\; V(\Delta\varphi)\cdot I(\Delta\varphi) \;=\; \mathcal{C}_3(\Delta\varphi)\,\mathcal{S}_3(\Delta\varphi) \;=\; (128 - \Delta\varphi)\,\Delta\varphi \;=\; 64^2 - (64 - \Delta\varphi)^2 \;\le\; \mathbf{4\,096 \;=\; K_P(3)^2},
  \end{equation}
  достигая абсолютного максимума ровно на биссектрисе октанта $\Delta\varphi^* = K_P(3) = 64$ ($V = I = 64$, закон согласования импедансов). При этом, пока пошаговый сдвиг фазы между соседними узлами решётки удерживается внутри бесфононного коридора $|\Delta\varphi| \le 40 < 41$, число рождаемых тепловых фононов тождественно равно $\hat{n}_E(\Delta\varphi) = \lfloor|\Delta\varphi|/41\rfloor = \mathbf{0}$ --- фазовая эстафета бежит по проводнику \textbf{без омического сопротивления и без выделения тепла (сверхток)}! А омическое сопротивление и джоулев нагрев возникают тогда и только тогда, когда на примесях, дефектах или тепловых колебаниях решётки локальный скачок фазы пробивает порог дискретности $|\Delta\varphi| \ge 41$, сбрасывая $\hat{n}_E(\Delta\varphi) \in \{1,\dots,6\}$ квантов энергии в фононный нагрев провода (и одновременно рождая вокруг провода магнитные вихревые плакеты $B_f = 512\,n_f$ по формуле~\eqref{eq:course-stokes-chern} из-за скачка фазы между проводом и окружающим воздухом).
\item \textbf{Полная теоретическая электротехника (ТОЭ) на кольце $\mathbb{Z}_{512}$: постоянный и переменный ток, $RLC$-цепи, выпрямители и трансформаторы:}
  \begin{itemize}[noitemsep]
  \item \emph{Постоянный (DC) и переменный (AC) ток:} постоянный ток есть стационарный пространственный уклон фазы $\Delta\varphi_{\mathrm{DC}} = k_{\mathrm{DC}} = \mathrm{const}\in[0,128]$ вдоль цепи ($V_{\mathrm{DC}} = 128 - k_{\mathrm{DC}},\ I_{\mathrm{DC}} = k_{\mathrm{DC}}$), а переменный ток --- вращение самой фазы $\varphi(t)\in\mathbb{Z}_{512}$ по всем четырём квадрантам Грея с нулевым средним за период ($\sum_{\varphi=0}^{511}\mathcal{C}_3(\varphi) = \sum_{\varphi=0}^{511}\mathcal{S}_3(\varphi) = 0$) и активной энергией Парсеваля $\sum_{\varphi=0}^{511}\mathcal{C}_3(\varphi)^2 = \sum_{\varphi=0}^{511}\mathcal{S}_3(\varphi)^2 = \mathbf{2\,796\,288}$.
  \item \emph{Законы Кирхгофа (KVL и KCL):} второй закон Кирхгофа для напряжений в замкнутом контуре есть кограничное тождество $(d_1 d_0\varphi)_f \equiv 0$ (в отсутствие переменного магнитного потока) и закон индукции Фарадея $(d_1 r)_f = 512\,n_f$ (при захвате вихревых плакет), а первый закон Кирхгофа для токов в узле следует из $d_2\circ d_1 \equiv 0$ и антиподального баланса 12 лучей кубооктаэдра $\sum_{k=0}^{11}\vec{e}_k = \mathbf{0}$.
  \item \emph{Реактивные элементы $L, C$ и резонанс $LC$-контура:} сдвиг фазы на $90^\circ$ на кольце $\mathbb{Z}_{512}$ есть переключение ровно на одну четверть кольца $Q_\square(3) = 128$ тиков ($\mathcal{C}_3(\varphi+128) = -\mathcal{S}_3(\varphi),\ \mathcal{S}_3(\varphi+128) = +\mathcal{C}_3(\varphi)$). Реактивная мощность за период тождественно равна нулю в силу ортогональности чётного и нечётного каналов ротатора:
    \[
      \sum_{\varphi=0}^{511} \mathcal{C}_3(\varphi)\,\mathcal{S}_3(\varphi) \;\equiv\; \mathbf{0},
    \]
    а в последовательном $LC$-контуре на резонансной частоте напряжения на индуктивности ($\varphi + 128$) и на ёмкости ($\varphi - 128 \equiv \varphi + 384$) отстоят друг от друга на полупериод противофазы $256$ тиков и взаимно уничтожаются в каждой из $512$ точек кольца: $\mathcal{C}_3(\varphi+128) + \mathcal{C}_3(\varphi+384) \equiv \mathbf{0}$!
  \item \emph{Однополупериодный, мостовой (Гретца) и беспульсационный многофазный выпрямители:} безветвевой однополупериодный диодный выпрямитель $\mathrm{rect}_{1/2}(x) = (|x|+x)/2$ пропускает положительный полупериод ($b_8 = 0$), накапливая за период заряд $\sum_{\varphi=0}^{511}\mathrm{rect}_{1/2}(\mathcal{S}_3(\varphi)) = 128^2 = \mathbf{16\,384}$ со средним выпрямленным током $16\,384 / 512 = \mathbf{32 = K_P(3)/2}$. Двухполупериодный диодный мост Гретца из $4$ диодов инвертирует знаковый бит Грея $b_8$, давая за период заряд $\sum_{\varphi=0}^{511}|\mathcal{S}_3(\varphi)| = 2\times 128^2 = \mathbf{32\,768 = 4\varepsilon(3)}$ и точный рациональный средний ток $32\,768 / 512 = \mathbf{64 = K_P(3) = I_{\max}/2}$. В многофазных системах переменного тока (с любым числом фаз $M\ge 1$) двухфазный квадратурный мостовой выпрямитель ($\varphi$ и $\varphi - 128$) в силу $L^1$-инварианта ротатора даёт \textbf{абсолютно постоянное выпрямленное напряжение с нулевыми пульсациями ($0\%$ ripple) без сглаживающего конденсатора}:
    \[
      |\mathcal{C}_3(\varphi)| \;+\; |\mathcal{S}_3(\varphi)| \;\equiv\; \mathbf{128 \;=\; Q_\square(3)} \qquad \forall\,\varphi\in\mathbb{Z}_{512}!
    \]
  \item \emph{Трансформатор и согласование импедансов витками $N_1 : N_2$:} поскольку общий замкнутый магнитопровод первичной ($N_1$ витков) и вторичной ($N_2$ витков) обмоток пронизывается единым квантованным потоком плакет $\Phi_{\mathrm{core}} = 512\,n_f$ ($d_2\circ d_1 \equiv 0$), на каждом витке индуцируется одинаковая ЭДС $\varepsilon_1 = \Delta_t\Phi_{\mathrm{core}}$, откуда $V_1 = N_1\varepsilon_1$, $V_2 = N_2\varepsilon_1$ ($N_2 V_1 - N_1 V_2 \equiv 0$). При балансе ампер-витков $N_1 I_1 - N_2 I_2 \equiv 0$ ($I_1 = N_2 i_0,\ I_2 = N_1 i_0$) мощность сохраняется тождественно ($V_1 I_1 = N_1 N_2 \varepsilon_1 i_0 \equiv V_2 I_2$), а входной импеданс преобразуется квадратично по коэффициенту трансформации: $Z_{\mathrm{in}} = (N_1 / N_2)^2 Z_{\mathrm{load}}$.
  \end{itemize}
\end{enumerate}

\section{Урок 11. Почему невозможен Демон Лапласа (\texorpdfstring{$64 \times 81 = 72^2$}{64 x 81 = 72\^2}), но работает 12-лучевой Геометрический Демон Максвелла}

В 1814~году Пьер-Симон Лаплас предположил: если закон движения Вселенной детерминирован, то достаточно мощный ум («Демон Лапласа»), измерив точное состояние всех частиц, сможет вычислить всё будущее. Чтобы спастись от Демона Лапласа, физика XX~века отказалась от микродетерминизма.

В нашей теории микрозакон Верле \textbf{на 100\% детерминирован и побитово обратим}, однако \textbf{Демон Лапласа внутри решётки математически невозможен} сразу по трём причинам:
\begin{enumerate}[label=\textbf{\arabic*.},noitemsep]
\item \textbf{Запрет самовложенности регистра Дирихле:} демон внутри решётки занимает собственную подсистему $V_{\mathrm{demon}} < \mathcal{V}_{O_2}(3,T)$ узлов причинного конуса, поэтому его память $72\,V_{\mathrm{demon}}$ бит строго меньше фазового объёма конуса $72\,\mathcal{V}_{O_2}(3,T)$ бит ($72 < 936$ уже при $T=1$ и $936 < 10\,584$ при $T=3$).
\item \textbf{Точный квадратичный закон слепоты Т-прибора ($64\times 81 = 72^2 = 5\,184$):}
  любой прибор взаимодействует с ячейкой через ворота Т-уровня, где расслоение Вейля--Швингера скрывает $K_P(3) = \mathbf{64}$ состояний мантиссы (Уровень~I), а мёртвая зона затвора $|\Delta\varphi|\le 40$ скрывает $2\Delta\varphi_{\mathrm{disc}}(3) - 1 = 2\times 41 - 1 = \mathbf{81 = 3^4}$ подпороговых состояний с нулевым излучением $\hat{n}_E = 0$ (Уровень~II). Их произведение в точности равно \textbf{квадрату двухтактового битового бюджета ячейки}:
  \begin{equation}
    \label{eq:course-laplace-5184}
    K_P(3)\cdot |\mathcal{E}_{\mathrm{Heis}}|(3)
    \;=\; 64 \times 81
    \;=\; 8^2 \times 9^2
    \;=\; (8\times 9)^2
    \;=\; \mathbf{72^2 \;=\; \bigl(B_2(3)\bigr)^2 \;=\; 5\,184}!
  \end{equation}
  Пассивный наблюдатель на Т-уровне не различает $5\,184$ микросостояния в каждой ячейке, а если он пошлёт надпороговый зонд $|\Phi_{\mathrm{probe}}|\ge 41$ ($\hat{n}_E\ge 1$), он сам же сдвинет шаг Верле измеряемой ячейки и изменит её будущее!
\item \textbf{Горизонт скорости звена $c_0 = \sqrt{2}\,c$:} данные со сферы $S_T$ обратного конуса $\mathcal{V}_{O_2}(3,T)$ идут к демону со скоростью 1 оболочка за такт и приходят в центр ровно в такт~$T$, когда само событие уже свершилось!
\end{enumerate}

\noindent\textbf{А что же с Демоном Максвелла?} Классический демон Сциларда--Бриллюэна измерял молекулы, записывал биты в память и платил тепло Ландауэра за стирание этой памяти. Но наш 12-лучевой кубооктаэдр реализует \textbf{безпамятьевого Геометрического Демона Максвелла}: так как $12\times 41 = 492 > 256 = \Phi_{\mathrm{Pauli}}$, даже максимальный 6-квантовый тепловой удар $\Phi_{\mathrm{Pauli}} = 256$ ($\hat{n}_E(256)=6$), распределённый по $12$ лучам кубооктаэдра, даёт на каждом луче:
\[
  \lfloor 256 / 12 \rfloor \;=\; \mathbf{21 \;<\; 41 \;=\; \Delta\varphi_{\mathrm{disc}}(3)}
  \qquad\Longrightarrow\qquad
  \hat{n}_E(21) \;=\; \mathbf{0\text{ тепловых фононов}},
\]
выдавая при этом строго положительный когерентный выпрямленный ток $12 \times \mathcal{S}_3(21) = 12 \times 21 = \mathbf{252 > 0}$ без записи и без стирания единого бита памяти!

